\documentclass[11pt]{article}

% ---- Times New Roman look, standard margins, single column (article default) ----
% Tightened deliberately to fit on ONE page -- an instruction sheet a clinician skims should not
% spill to a second page with a few lines on it.
\usepackage[margin=0.85in,top=0.55in,bottom=0.55in]{geometry}
\usepackage{newtxtext,newtxmath}   % Times New Roman-equivalent text + math
\usepackage[T1]{fontenc}
\usepackage{booktabs}
\usepackage{array}
\usepackage{enumitem}
\usepackage{parskip}
\usepackage{titlesec}
\usepackage[hidelinks]{hyperref}

\setlength{\parskip}{3pt}
\titlespacing*{\section}{0pt}{7pt}{3pt}
\titleformat{\section}{\large\bfseries}{\thesection}{0.6em}{}
\titleformat{\subsection}{\normalsize\bfseries}{\thesubsection}{0.6em}{}
\titlespacing*{\subsection}{0pt}{6pt}{2pt}
\setlist[enumerate]{leftmargin=1.6em, itemsep=1pt, topsep=1pt}
\setlist[itemize]{leftmargin=1.6em, itemsep=1pt, topsep=1pt}

\newcolumntype{L}[1]{>{\raggedright\arraybackslash}p{#1}}

\pagenumbering{gobble}

\begin{document}

\begin{center}
{\Large\bfseries Respiratory Sound Listening Task}\\[2pt]
{\normalsize Instructions}
\end{center}

\vspace{2pt}
\noindent Thank you for helping with this. It should take about \textbf{75 minutes}, and you can
stop and resume at any time --- nothing needs to be finished in one sitting.

\section*{Setup (2 minutes)}
\begin{enumerate}
  \item Use \textbf{headphones or earbuds}, not laptop speakers --- the sounds we are asking about
        are quiet and very brief, and laptop speakers cannot reproduce them clearly.
  \item Sit somewhere quiet. Set a comfortable volume using
        \texttt{reference\_sounds/example\_normal.wav}, then \textbf{do not change the volume}
        afterward --- every clip is recorded at the same loudness.
\end{enumerate}

\section*{Step 1: Calibrate (5 minutes)}
Listen to every file in the \texttt{reference\_sounds/} folder before labeling anything. These
examples show what we mean by each term below, so we are both using the same words for the same
sounds. Return to them at any time during the task if you are unsure.

\emph{Note: these reference examples are drawn from the dataset's own existing labels, not
selected by a clinician. If one sounds mislabeled to you, please tell us --- that is useful
information on its own.}

\section*{Step 2: Label the clips (about 70 minutes)}

Open \texttt{labels.xlsx}. Every answer cell is a dropdown menu, so you select an option rather
than type it.

\begin{enumerate}
  \item Play \texttt{clips/clip\_001.wav} in any media player (double-clicking usually works).
  \item Fill in the corresponding row using the dropdown menus. Replay as many times as you like
        before answering.
  \item Move to the next clip. \textbf{Please go in order and do not skip ahead.} If a clip is too
        unclear to judge, choose \texttt{unsure} and move on --- do not guess.
\end{enumerate}

\subsection*{The columns}

\begin{center}
\renewcommand{\arraystretch}{1.15}
\begin{tabular}{@{}L{2.7cm}L{5.6cm}L{5.5cm}@{}}
\toprule
\textbf{Column} & \textbf{Options} & \textbf{Notes} \\
\midrule
\texttt{crackles} & none / fine / coarse / both / unsure & The one we care about most \\
\texttt{wheeze} & none / wheeze / rhonchi / both / unsure & \\
\texttt{confidence} & low / medium / high & Please be honest --- low-confidence answers are
   still useful \\
\texttt{audio\_quality} & ok / noisy / unusable & Lets us exclude a bad recording rather than
   blame a good detector for it \\
\texttt{notes} & free text (optional) & Anything else worth telling us \\
\bottomrule
\end{tabular}
\end{center}

\noindent\textbf{Please use \texttt{unsure} freely.} An honest ``unsure'' is far more useful to us
than a guess. If you mark many clips \texttt{unsure}, that is a genuine and useful finding for
us --- it will not mean your time was wasted.

\section*{What we are \emph{not} asking}
You are not diagnosing a patient. There is no clinical history, no imaging, and nothing here
affects anyone's care. These are anonymous archived research recordings. The only question at any
point is simply: \emph{what sound do you hear in this clip?}

\section*{When you are finished}
If you were sent a \textbf{Google Sheets} link, your answers are already saved --- just let us
know you are done. Otherwise, please email the completed \texttt{labels.xlsx} file back to us.

\vspace{6pt}
\noindent Thank you again for your time and expertise --- this is genuinely helpful to the
project.

\end{document}
